\newcommand{\R}{\mathbb{R}}
\newcommand{\Borel}[1]{\mathscr{B}(\mathbb{R}^{#1})}
\newcommand{\bors}{\mathscr{B}}
\newcommand{\simp}{\mathscr{S}}
\newcommand{\E}{\mathbb{E}}
\appendix

\section{Fundamentals}

\subsection{Wiener Integral}

Let $T$ be a set and $X := \{X(t)\}_{t\in T}$ a $T$-indexed stochastic process. We recall that $X$ is a Gaussian random field (process when $T \subset \mathbb{R}$) if $(X_{t_1} ,\cdots,X_{t_m})$ is a Gaussian random vector for all $t_1,\cdots,t_m \in T$.\par

\begin{definition}
Let $\Borel{m}$ denote the collection of all Borel-measurable subsets of $\mathbb{R}^m$ that have finite Lebesgue measure. White noise on $\mathbb{R}^m$ is a mean-zero set-indexed Gaussian random field ${\xi(A)}_{A\in \Borel{m}}$ with covariance function
\[E[\xi(A_1)\xi(A_2)] := |A_1 \cap A_2| \quad\text{ for all } A_1, A_2 \in \mathscr{L}(\mathbb{R}^m), \]
where $|\cdot|$ denotes the Lebesgue measure on $\mathbb{R}^m$ for every $m$.
\end{definition}
\emph{Remark.} Here is a general proof from Ash and Gardner.
\begin{proof}
    For any complex numbers $c_1,\cdots,c_n$ and $A_1,\cdots,A_n \in \Borel{m}$, we have
    \[
    \sum\limits_{1\leq i,j\leq n} c_i\bar{c_j} |A_i\cap A_j| = \sum\limits_{1\leq i,j\leq n}c_i\bar{c_j}\int \chi(A_i)\chi(A_j) = \int \left|\sum\limits_{i=1}^n c_i\chi(A_i)\right|^2 \geq 0
    \]
    and we are done.
\end{proof}

\begin{lemma}
    $\{\xi(A)\}_{A \in \Borel{m}}$ is an $L^2(\Omega)$-valued countably-additive measure on $\Borel{m}$.
\end{lemma}

\begin{definition}[Wiener Integral]\par
    For $h \in \simp(\R^m)$ the set of the simple functions, define
    \[
    \int h d\xi := \sum\limits_{i=1}^n c_i\xi(A_i)
    \]
    where $h(x) = \sum\limits_{i=1}^n c_i\chi(A_i)$ for some $A_1,\cdots,A_n\in\bors(\R^m)$ and $c_1,\cdots,c_n \in \R$.\par
    For $f\in L^2(\R^m)$, define 
    \[
    \int f d\xi = \lim_{n\to\infty} \int f_i d\xi
    \]
    in $L^2(\Omega)$ for some $f_i\in \simp(\R^m)$ and $f_i \to f$ in $L^2$.
\end{definition}
\begin{proof}
    
\end{proof}

\begin{proposition}
    The $L^2(\R^m)$-indexed Gaussian random field $\{\int h d\xi\}_{h\in L^2(\R^m)}$ has the following properties:
    \begin{enumerate}
        \item $\int (a_1h_1+a_2h_2)d\xi = a_1\int h_1d\xi + a_2 \int h_2 d\xi$ a.s. for all $a_1,a_2 \in \R$ and $h_1,h_2 \in L^2(\R^m)$.
        \item For all $h, h_1,h_2 \in L^2(\R^m)$,
        \[\E \int hd\xi = 0, \E(\int h_1d\xi \int h_2d\xi) = \int h_1h_2\]
    \end{enumerate}
\end{proposition}
    \emph{Remark.} We may define
    \[\int_A h d\xi = \int h\chi(A)d\xi\]
    and it is easy to prove the \emph{Wiener isometry} by the properties above.

\begin{definition}
    For every $f\in L^2(\R^m)$ and $x\in\R^m$, define
    \[
    (f*\xi)(x) := \int f(x-y)\xi(dy)
    \]
\end{definition}

\begin{proposition}
    For every $f\in L^2(\R^m)$, there exists a version of $f*\xi$ that is measurable jointly in $(\omega,x)$. Moreover, we have the following Young-type inequality: For every finite Borel measure $\mu$ on $\R^m$,
    \[
    \E\left(\left|\int|(f*\xi)(x)|\mu(dx)\right|^2\right) \leq \mu(\R^m)^2||f||^2_{L^2(\R^m)}
    \]
\end{proposition}
\begin{proof}
    
\end{proof}

\subsection{Useful Approximations}

\begin{proposition}[Bihari-LaSalle Inequality]\ \par
    Let $u$ and $f$ be non-negative continuous functions defined on the half-infinite ray $[0,\infty)$ and let $w$ be a continuous non-decreasing function defined on $[0,\infty)$ and $w(u) > 0$ on $(0,\infty)$. For $u$ we have
    \[
    u(t) \leq \alpha + \int_0^t f(s)w(u(s)) ds,\quad t\in[0,\infty),\label{BLIneq}
    \]
    where $\alpha$ is a non-negative constant, then
    \[u(t) \leq G^{-1}\left(G(\alpha)+\int_0^t f(s) ds\right),\quad t\in[0,T]\]
    where the function $G$ is defined by
    \[
    G(x) = \int_{x_0}^x \dfrac{dy}{w(y)},\quad x\geq 0, x_0 \geq 0,
    \]
    and $G^{-1}$ is the inverse function of $G$ and $T$ is chosen so that
    \[
    G(\alpha) + \int_0^t f(s)ds\in\text{Dom}(G^{-1}),\quad \forall t\in[0,T]
    \]
\end{proposition}
\begin{proof}
    Notice $G(x) = \int_{x_0}^x \dfrac{dy}{w(y)}$ is increasing and
    \[
    \text{Dom}(G^{-1}) = \left[G(0)=-\int_{0}^{x_0}\dfrac{dy}{w(y)}, G(0) + \int_0^{\infty}\dfrac{dy}{w(y)}\right]
    \]
    then since $u\geq 0$, then we have
    \[
    G(u(t)) \leq G\left(\alpha + \int_0^t f(s) \omega(u(s))ds\right):=H(t)
    \]
    and we have
    \[
    H'(t) = \dfrac{f(t)\omega(u(t))}{\omega(\alpha+\int_0^t f(s)\omega(u(s))ds)} \leq f(t)
    \]
    so we have
    \[
    H(t) \leq G(\alpha) + \int_0^t f(s)ds.
    \]
    Therefore, we may have the requested inequality for all $0\leq t \leq T$ where
    \[
    T := \sup_{t}\left\{\int_0^tf(s)ds \leq G(0)-G(\alpha)+\int_0^{\infty}\dfrac{dy}{\omega(y)}\right\}
    \]
\end{proof}

\subsection{Matrix Analysis}

\begin{theorem}[Spectral Theorem for Real Symmetric Matrices]\ \par
Let $A \in \mathbb{R}^{n\times n}$ be a real symmetric matrix, that is, $A^T = A$. Then:

\begin{enumerate}
\item All eigenvalues of $A$ are real.
\item Eigenspaces corresponding to distinct eigenvalues are orthogonal.
\item The space $\mathbb{R}^n$ decomposes as an orthogonal direct sum of eigenspaces:
\[
\mathbb{R}^n = \bigoplus_{\lambda \in \sigma(A)} \ker(A - \lambda I).
\]
\item Equivalently, there exists an orthogonal matrix $Q$ and a real diagonal matrix $\Lambda$ such that
\[
A = Q \Lambda Q^T.
\]
\end{enumerate}
\end{theorem}

\subsection{Kolmogorov's Theorems}

\begin{definition}
    For any $w\in\R^m$, then
    \[\rho(w):=\sum\limits_{i=1}^m |w_j|^{\alpha_j}\]
    where $0 < \alpha_1,\cdots,\alpha_m \leq 1$ are fixed.\par
    For a bounded and closed subset $T$ of $\R^m$, define
    \[
    \Lambda_T(a):= \int_T dt\int_T ds[\rho(t-s)]^{-a}\quad\text{ for all }a>0.
    \]
\end{definition}
\emph{Remark. }$\Lambda_T(a) <\infty$ if and only if $a< H:= \sum\limits_{j=1}^m \alpha_j^{-1}$.

\begin{theorem}[Kolmogorov's Continuity Theorem]\ \par
    Let $\{X_t\}_{t\in T}$ be a stochastic process, where $T\subset \R^m$ is measurable and bounded, and suppose that there exists finite real numbers $C>0$ and $k>H$, such that
    \[\lVert X_t - X_s\rVert_k \leq C\rho(t-s)\quad\text{ for any }s,t\in T\]
    Then $X$ has a modification $\overline{X}$ that is Holder continuous a.s. Moreover, $\overline{X}$ satisfies the following: For all $q\in (0,1-(H/k))$ and $\delta \in (q,1-(H/k))$ there exists a finite constant $D$, that does not depend on that numerical value of $k$, such that
    \[
    \E\left(\sup_{s,t\in T,s\neq t}\left|\dfrac{\overline{X}_t-\overline{X}_s}{[\rho(t-s)]^q}\right|^k \right) \leq (\delta -q)^{-1}\left(\dfrac{DC}{\delta}\right)^k k^{2H}\Lambda_T(2H-k+k\delta)
    \]
\end{theorem}